% !TEX root = ../discretemath.tex

\section{Обобщённая теорема Ван дер Вардена. Монохромный квадрат в $\mathbb Z^2$}


\subsection{Формулировка обобщённой теоремы}

\begin{Theorem}{Обобщённая теорема Ван дер Вардена}{gvdw}
	Пусть $n,r\in\mathbb N$ и $G\subset\mathbb Z^n$ — конечное множество. При любой $r$-раскраске $\mathbb Z^n$ найдётся одноцветная \emph{гомотетичная копия} $G$, то есть множество
	\[
	a+\lambda G=\{a+\lambda g\mid g\in G\},\qquad a\in\mathbb Z^n,\ \lambda\in\mathbb N,
	\]
	все точки которого покрашены в один цвет.
\end{Theorem}

\begin{Remark}{}{}
	При $n=1$, $G=\{0,1,\ldots,k-1\}$ копия $a+\lambda G=\{a,a+\lambda,\ldots,a+(k-1)\lambda\}$ — арифметическая прогрессия длины $k$, и теорема превращается в обычную теорему Ван дер Вардена. Общий случай доказывается индукцией по размерности $n$, по $|G|$ и по числу цветов $r$; на каждом шаге применяется одномерная теорема.
\end{Remark}

\subsection{Монохромный квадрат в $\mathbb Z^2$}

\begin{Statement}{}{}
	При любой $2$-раскраске $\mathbb Z^2$ найдётся монохромный ось-выровненный квадрат — четвёрка точек
	\[
	(x,y),\ (x{+}d,y),\ (x,y{+}d),\ (x{+}d,y{+}d),\qquad d>0,
	\]
	одного цвета.
\end{Statement}

Это частный случай обобщённой теоремы для $G=\{(0,0),(1,0),(0,1),(1,1)\}$. Ниже — прямое доказательство блочным аргументом (оно же иллюстрирует общую схему: «поднятие» одномерной теоремы за счёт укрупнения алфавита цветов).

\begin{proof}
	Пусть $N_2$ — параметр, который выберем позже.

	\textbf{Шаг 1 (блоки).} Разобьём плоскость на вертикальные полосы ширины $N_2$. Раскраска блока $N_2\times N_2$ — один из $2^{N_2^2}$ вариантов (его «блочный цвет»).

	\begin{center}
		\begin{tikzpicture}[scale=0.7]
			\foreach \x in {0,2,4,6,8}{
				\draw[gray!40,fill=gray!8] (\x,0) rectangle (\x+1.7,2.5);
			}
			\draw[blue!60,fill=blue!15,thick] (0,0) rectangle (1.7,2.5);
			\draw[blue!60,fill=blue!15,thick] (4,0) rectangle (5.7,2.5);
			\node[font=\small] at (0.85,-0.4) {$B_1$};
			\node[font=\small] at (4.85,-0.4) {$B_2$};
			\draw[<->,thick,blue!70] (0,-0.85) -- (4,-0.85)
			node[midway,below,font=\footnotesize]{расстояние $D$};
			\node[font=\footnotesize] at (10.2,1.25) {$\cdots$};
			\draw[<->,thin] (-0.2,0) -- (-0.2,2.5) node[midway,left,font=\footnotesize]{$N_2$};
		\end{tikzpicture}
	\end{center}

	\textbf{Шаг 2 (три одинаковых блока).} Рассмотрим горизонтальный ряд блоков $B_1,B_2,\ldots$ как одномерную раскраску в $2^{N_2^2}$ «блочных цветов». По одномерной теореме Ван дер Вардена (для длины $3$) при достаточном числе блоков найдутся \emph{три} идентично раскрашенных блока $B,B{+}D,B{+}2D$ в арифметической прогрессии.

	\textbf{Шаг 3 (нахождение квадрата).} Три идентичных блока задают одинаковое правило раскраски в каждой строке. В пределах блока высоты $N_2$ рассмотрим пары строк $(r,r{+}D)$ с $r{+}D<N_2$. Внутри фиксированного столбца $x_0$ цвета строк $r$ и $r{+}D$ — лишь $2$ варианта на каждую из $N_2$ строк, поэтому при $N_2>2D$ найдётся пара строк $r,r{+}D$ одного цвета $c$ (принцип Дирихле по парам строк). Поскольку блоки $B$ и $B{+}D$ (отстоящие на $D$ по горизонтали) раскрашены одинаково, в обоих столбцах $x_0$ и $x_0{+}D$ строки $r$ и $r{+}D$ имеют тот же цвет $c$. Тогда четвёрка
	\[
	(x_0,r),\ (x_0{+}D,r),\ (x_0,r{+}D),\ (x_0{+}D,r{+}D)
	\]
	— монохромный квадрат со стороной $D$.
\end{proof}

\begin{figure}[H]
	\centering
	\begin{tikzpicture}[scale=0.62]
		\foreach \x in {0,...,7}
		\foreach \y in {0,...,5}
		\fill[gray!35] (\x,\y) circle (2.2pt);
		\draw[red!80!black, very thick] (1,0) -- (5,0) -- (5,4) -- (1,4) -- cycle;
		\foreach \p in {(1,0),(5,0),(5,4),(1,4)}{ \fill[red!75!black] \p circle (4.5pt); }
		\node[below=5pt, font=\small, red!80!black] at (3,-0.1) {монохромный квадрат со стороной $D$};
	\end{tikzpicture}
	\caption{Итог блочного аргумента: четыре вершины (красные) одного цвета образуют ось-выровненный квадрат со стороной $D$ — гомотетичную копию $G=\{(0,0),(1,0),(0,1),(1,1)\}$.}
\end{figure}
